National government & Political/reputational risk if the first project is delayed, over budget or meets public resistance; fiscal risk if a subsidy (CfD/SDE-type) is required. & Contribution to climate and renewable targets, energy security and less import dependence; start of a national offshore wind industry and jobs. \\
\hline
Local government & Visual impact on the coastline and possible effect on tourism and property values; nuisance from onshore cable route and substation works. & Local jobs, business rates/taxes, community fund, upgraded port and grid infrastructure, image as a `green' region. \\
\hline
Nature conservation society & Bird and bat collisions, underwater noise during pile driving (marine mammals), habitat disturbance from cables and scour protection. & Artificial reef effect and nature-inclusive design, no-trawling zone as de-facto marine protected area, monitoring data, displacement of fossil generation. \\
\hline
Local fishermen & Loss of fishing grounds and longer sailing routes; risk of gear snagging on cables; reduced catch during construction. & Compensation, guard-vessel and survey contracts during construction and O\&M, passive fishing/aquaculture within the park, higher fish stock near foundations. \\
\hline
Local port & Congestion and conflicts with existing users; investment in quays/storage that may not pay back if no follow-up projects come. & Marshalling and O\&M base for decades: new revenue, jobs, positioning as offshore wind hub for future projects. \\
\hline
Grid company (TSO) & Grid congestion and balancing effort due to variable output; required onshore reinforcements; technical risk of a first offshore connection (reactive power, harmonics). & Large new generation close to the coast, learning on offshore grid standards; possible reuse of offshore substation for future grid hub/interconnector. \\
